\documentclass[11pt,a4paper]{article}

\usepackage[utf8]{inputenc}
\usepackage[T1]{fontenc}
\usepackage[english]{babel}
\usepackage{amsmath,amssymb,amsthm}
\usepackage{graphicx}
\usepackage{booktabs}
\usepackage{hyperref}
\usepackage{geometry}
\usepackage{algorithm}
\usepackage{algorithmic}
\geometry{margin=2.5cm}

\title{Comparing optimization algorithms for image classification on CIFAR-10:\\
       a convex / non-convex perspective}
\author{Ilyan \textsc{Debbah} \and Alexan \textsc{Lastname} \and Gabriel \textsc{Lastname} \and Thi-Tho \textsc{Lastname}\\
        \small EFREI — Convex optimization, M.~B\textsc{enedetti}, 2026}
\date{June 2026}

\begin{document}
\maketitle

\begin{abstract}
We study the empirical behaviour of several first- and second-order optimization
algorithms (gradient descent, SGD with momentum, Nesterov, AdaGrad, RMSProp,
Adam, L-BFGS) on the CIFAR-10 image classification task. To bridge the gap
between the convex theory developed in the course and the predominantly
non-convex losses used in deep learning, we run the same optimizers in two
regimes: (i)~a strictly convex one --- multinomial logistic regression and
linear SVM on raw pixels --- where the rates predicted by the theory hold, and
(ii)~a non-convex one --- a small convolutional network --- where they do not.
The comparison highlights where the convex theory transfers and where it
breaks down.
\end{abstract}

\tableofcontents

% =============================================================================
\section{Introduction}
% =============================================================================
% TODO:
% - Motivate: ML losses are usually non-convex, yet we train them with
%   algorithms designed for the convex case. Why does it work?
% - Briefly introduce the dataset (CIFAR-10) and why we picked it (richer
%   than MNIST, enough to stress optimizers).
% - State the contribution: a side-by-side comparison of the same optimizers
%   on a convex baseline and a non-convex CNN, on the same data.

% =============================================================================
\section{Background and problem formulation}
% =============================================================================

\subsection{Convex optimization in machine learning}
% TODO:
% - Recall the empirical risk minimization formulation.
% - Recall convex functions, smoothness L, strong convexity mu, condition kappa.
% - State the rates we will compare against: O(1/k) for GD on smooth convex,
%   O(1/k^2) for Nesterov, linear rate for strongly convex, quadratic local
%   rate for Newton.

\subsection{The CIFAR-10 dataset}
% TODO:
% - 60k 32x32 colour images, 10 classes, balanced.
% - Refer to fig:dataset_samples.

% =============================================================================
\section{Models}
% =============================================================================

\subsection{Convex models}
\subsubsection{Multinomial logistic regression}
% L(W) = -1/N sum log softmax_y(Wx) + l2/2 ||W||^2
% Smoothness, strong convexity from the L2 term.

\subsubsection{Linear SVM (squared hinge)}
% L(W) = 1/N sum_i sum_{j != y_i} max(0, s_j - s_y + 1)^2 + l2/2 ||W||^2

\subsection{Non-convex model: a small CNN}
% Architecture (see src/cnn_model.py); ~200k parameters.
% Non-convex because of (i) ReLU non-linearities, (ii) composition of layers.

% =============================================================================
\section{Optimization algorithms studied}
% =============================================================================
% For each algorithm: update rule, known convergence rate in the convex
% setting, intuition for behaviour on non-convex problems, hyperparameters.

\subsection{Gradient descent and stochastic gradient descent}
\subsection{Momentum and Nesterov accelerated gradient}
\subsection{Adaptive methods: AdaGrad, RMSProp, Adam}
\subsection{Quasi-Newton: L-BFGS}

% =============================================================================
\section{Experimental protocol}
% =============================================================================
% - Hardware, framework versions, seeds.
% - Train/test splits.
% - Hyperparameter search procedure (grid for lr, momentum).
% - Metrics: training loss, gradient norm, test accuracy, wall-clock time.

% =============================================================================
\section{Results}
% =============================================================================

\subsection{Convex regime}
% Figures: loss curves (log scale), gradient norm, test accuracy table.
% Comments: empirical rates vs theory, role of L2, conditioning.

\subsection{Non-convex regime}
% Same figures. Discussion of differences with the convex regime.
% Why does the gradient norm not go to zero?
% Why does Adam dominate early epochs?
% Generalization comments (SGD+momentum vs Adam).

\subsection{Cross-regime comparison}
% Side-by-side table summarising the ranking of optimizers in each regime.
% Discussion: where convex theory transfers, where it breaks down.

% =============================================================================
\section{Discussion: convex theory facing non-convex losses}
% =============================================================================
% Required by the project brief (section 1.2 of the assignment).
% - Why first-order methods still work on non-convex losses (overparametrization,
%   benign non-convexity, escape from saddle points).
% - Why 2nd-order methods do not scale to deep models (Hessian-vector cost,
%   non-positive-definite Hessian).
% - Mention recent work that brings some convergence guarantees back into the
%   non-convex setting (e.g. NTK regime, PL inequality).

% =============================================================================
\section{Conclusion}
% =============================================================================

\bibliographystyle{plain}
\bibliography{references}

\end{document}
